% ============================================================
% Academic Research Agent 详细设计说明书（迭代三）
% 编译方式：xelatex -> xelatex
% ============================================================
\documentclass[12pt,a4paper,onecolumn]{ctexart}

\usepackage[top=2.5cm,bottom=2.5cm,left=2.8cm,right=2.8cm]{geometry}
\usepackage{amsmath,amssymb}
\usepackage[table,xcdraw]{xcolor}
\usepackage{hyperref}
\usepackage{enumitem}
\usepackage{booktabs}
\usepackage{longtable}
\usepackage{array}
\usepackage{listings}
\usepackage{caption}
\usepackage{float}
\usepackage{tikz}
\usetikzlibrary{shapes.geometric, arrows.meta, positioning, fit, calc, backgrounds, chains}
\emergencystretch=2em

\hypersetup{colorlinks=true, linkcolor=blue!60!black, urlcolor=blue!60!black}
\setcounter{tocdepth}{2}
\newcommand{\resizetikz}[2][\textwidth]{\resizebox{#1}{!}{#2}}

\lstset{
  basicstyle=\small\ttfamily,
  keywordstyle=\bfseries\color{blue!60!black},
  commentstyle=\color{green!50!black},
  stringstyle=\color{red!50!black},
  frame=single,
  numbers=left,
  numberstyle=\tiny,
  breaklines=true,
  showstringspaces=false,
  tabsize=2
}

\title{\textbf{Academic Research Agent 详细设计说明书}}
\author{软件工程课程项目组}
\date{迭代三 \quad 2026年6月}

\begin{document}
\maketitle
\thispagestyle{empty}
\newpage
\tableofcontents
\newpage

%============================================================
\section{系统总体架构}
%============================================================

迭代三系统采用前端交互层、FastAPI 路由层、领域管理层、Agent Pipeline 层、文件系统持久化层的分层设计。相比迭代二的一次性执行，迭代三以 Session 作为核心数据边界，并将用户可编辑中间产物纳入 Pipeline。

\begin{figure}[H]
\centering
\resizetikz{\begin{tikzpicture}[
  node distance=0.9cm and 1.2cm,
  layer/.style={rectangle, draw=blue!65, fill=blue!6, rounded corners=2mm, text width=12.5cm, align=center, minimum height=0.9cm},
  comp/.style={rectangle, draw=teal!70, fill=teal!6, rounded corners=1mm, text width=3.2cm, align=center, minimum height=0.75cm},
  arr/.style={-{Stealth}, thick}
]
  \node[layer] (front) {前端交互层：home.html / index.html / notebooklm.js / Vditor Markdown 编辑器};
  \node[layer, below=of front] (routes) {FastAPI 路由层：pages / session / agent / chat / conversation / draft / admin};
  \node[layer, below=of routes] (manager) {领域管理层：SessionManager / SkillManager / GlobalKnowledgeBase / CopilotManager / ToolRegistry};
  \node[layer, below=of manager] (agent) {Agent Pipeline 层：Plan / ReAct Search / RAG Notes / Analysis / Writer / Skills};
  \node[layer, below=of agent] (storage) {持久化层：sessions/ 文件系统、config/tools.json、evaluation/ 评估数据};

  \draw[arr] (front) -- (routes);
  \draw[arr] (routes) -- (manager);
  \draw[arr] (manager) -- (agent);
  \draw[arr] (agent) -- (storage);
  \draw[arr] ([xshift=4.8cm]storage.north) .. controls +(1.2,1.2) and +(1.2,-1.2) .. ([xshift=4.8cm]manager.south);
\end{tikzpicture}}
\caption{迭代三系统分层架构}
\end{figure}

%============================================================
\section{后端模块设计}
%============================================================

\subsection{Web 入口}
\texttt{agent/web\_app.py} 作为轻量入口，负责创建 FastAPI 应用、初始化核心管理器、注入共享依赖并注册路由。业务逻辑被拆分到 \texttt{backend/routes/}，避免入口文件膨胀。

\subsection{路由模块职责}
\begin{longtable}{p{3.6cm}p{9.8cm}}
\toprule
\textbf{模块} & \textbf{职责} \\
\midrule
\texttt{deps.py} & 保存共享依赖、路径常量与运行时 \texttt{RUNS}。\\
\texttt{models.py} & 集中定义请求体 Pydantic 模型。\\
\texttt{pages.py} & 页面路由、收藏夹、历史记录与 PDF 文件响应。\\
\texttt{session.py} & Session、论文、笔记、分析卡片的持久化接口。\\
\texttt{agent.py} & Agent 阶段执行接口与自动模式后台线程。\\
\texttt{chat.py} & Session 内问答、RAG 上下文、修订判意与上下文压缩。\\
\texttt{conversation.py} & 多聊天会话管理。\\
\texttt{draft.py} & 综述草稿保存与读取。\\
\texttt{admin.py} & 工具管理、Skills、全局知识库与 Copilot。\\
\bottomrule
\end{longtable}

%============================================================
\section{Session 数据模型}
%============================================================

\subsection{目录结构}
\begin{lstlisting}[language=bash]
sessions/{session_id}/
+-- metadata.json
+-- plan/
|   +-- initial_plan.md
|   +-- confirmed_keywords.json
+-- papers/
|   +-- papers_list.json
|   +-- deleted_papers.json
|   +-- *.pdf / *.txt
+-- notes/
|   +-- draft_notes.md
+-- analysis/
|   +-- analysis_results.json
+-- draft/
|   +-- draft_v*.md
|   +-- current_draft.md
+-- traces/
|   +-- run_traces.json
+-- chats/
    +-- _index.json
    +-- conv_*.json
\end{lstlisting}

\subsection{状态机}
\begin{figure}[H]
\centering
\resizetikz{\begin{tikzpicture}[
  node distance=1.35cm,
  state/.style={rectangle, draw=purple!70, fill=purple!6, rounded corners=2mm, text width=2.6cm, align=center, minimum height=0.8cm},
  arr/.style={-{Stealth}, thick}
]
  \node[state] (planning) {planning};
  \node[state, right=of planning] (confirmed) {plan\\confirmed};
  \node[state, right=of confirmed] (searching) {searching};
  \node[state, right=of searching] (sc) {search\\complete};
  \node[state, below=1.1cm of sc] (notes) {reviewing\\notes};
  \node[state, left=of notes] (writing) {writing};
  \node[state, left=of writing] (draft) {reviewing\\draft};
  \node[state, left=of draft] (complete) {complete};

  \draw[arr] (planning) -- (confirmed);
  \draw[arr] (confirmed) -- (searching);
  \draw[arr] (searching) -- (sc);
  \draw[arr] (sc) -- (notes);
  \draw[arr] (notes) -- (writing);
  \draw[arr] (writing) -- (draft);
  \draw[arr] (draft) -- (complete);
  \draw[arr, dashed, gray] (complete.north) .. controls +(0,1.5) and +(0,1.5) .. node[above]{追加调研} (searching.north);
  \draw[arr, dashed, gray] (draft.north) .. controls +(0,1.0) and +(0,1.0) .. (searching.north);
\end{tikzpicture}}
\caption{Session 状态机}
\end{figure}

%============================================================
\section{Agent Pipeline 设计}
%============================================================

\subsection{分步模式与自动模式复用}
分步模式和自动模式复用同一组阶段函数。差异仅在控制权：分步模式由用户逐步触发，自动模式由后台线程按顺序执行。

\begin{figure}[H]
\centering
\resizetikz{\begin{tikzpicture}[
  node distance=1.4cm and 1.1cm,
  phase/.style={rectangle, draw=blue!70, fill=blue!6, rounded corners=2mm, text width=2.8cm, align=center, minimum height=0.85cm},
  human/.style={diamond, draw=orange!80, fill=orange!8, aspect=2.2, text width=2.4cm, align=center},
  arr/.style={-{Stealth}, thick}
]
  \node[phase] (plan) {run/plan};
  \node[human, right=of plan] (h1) {用户确认\\关键词};
  \node[phase, right=of h1] (search) {run/search};
  \node[human, right=of search] (h2) {用户筛选\\论文};
  \node[phase, below=1.2cm of h2] (notes) {run/notes};
  \node[human, left=of notes] (h3) {用户编辑\\笔记};
  \node[phase, left=of h3] (analysis) {run/analyze};
  \node[phase, left=of analysis] (write) {run/write};

  \draw[arr] (plan) -- (h1);
  \draw[arr] (h1) -- (search);
  \draw[arr] (search) -- (h2);
  \draw[arr] (h2) -- (notes);
  \draw[arr] (notes) -- (h3);
  \draw[arr] (h3) -- (analysis);
  \draw[arr] (analysis) -- (write);
  \draw[arr, dashed, gray] (plan.south) .. controls +(0,-2.2) and +(0,-2.2) .. node[below]{自动模式跳过人工暂停} (write.south);
\end{tikzpicture}}
\caption{分步模式与自动模式的控制差异}
\end{figure}

\subsection{ReAct 搜索循环}
搜索阶段由 \texttt{BaseAgent} 执行。每轮要求模型输出 JSON 格式的 \texttt{thought}、\texttt{action}、\texttt{action\_input}，工具返回写入 \texttt{observation}。

\begin{lstlisting}[language=Python]
while loop_count < max_loops:
    model_output = llm.chat(system_prompt, user_prompt, history)
    action = parse_json(model_output)
    observation = tool.execute(**action_input)
    traces.append({
        "thought": thought,
        "action": action_name,
        "observation": observation,
        "error_type": error_type,
    })
\end{lstlisting}

\subsection{错误反馈与质量门禁}
\begin{longtable}{p{3.1cm}p{5.4cm}p{4.7cm}}
\toprule
\textbf{机制} & \textbf{触发条件} & \textbf{处理方式} \\
\midrule
即时反馈 & JSON 解析错误、未知工具、工具异常 & 将错误写回下一轮上下文。\\
深度反思 & 连续同类错误达到阈值 & 要求 Agent 换策略。\\
FINISH 门禁 & 论文数量不足或笔记不足 & 阻止结束，提示继续检索。\\
PDF 兜底 & 搜索结束后 & 扫描 trace 中 arXiv ID 并尝试下载。\\
\bottomrule
\end{longtable}

%============================================================
\section{笔记、分析与写作模块}
%============================================================

\subsection{RAGNoteGenerator}
笔记生成模块优先使用 PDF 全文切块和 Embedding 检索。Embedding 失败时回退 BM25，PDF 缺失时回退摘要。

\begin{figure}[H]
\centering
\resizetikz{\begin{tikzpicture}[
  node distance=1.2cm,
  stepbox/.style={rectangle, draw=teal!70, fill=teal!6, rounded corners=2mm, text width=3.0cm, align=center},
  fallback/.style={rectangle, draw=red!60, fill=red!6, rounded corners=2mm, text width=3.2cm, align=center},
  arr/.style={-{Stealth}, thick}
]
  \node[stepbox] (pdf) {PDF 文本提取};
  \node[stepbox, right=of pdf] (chunk) {滑动窗口切块};
  \node[stepbox, right=of chunk] (emb) {Embedding/BM25 检索};
  \node[stepbox, right=of emb] (llm) {LLM 生成结构化笔记};
  \node[fallback, below=1.0cm of emb] (fb) {失败回退：摘要生成 / BM25};
  \draw[arr] (pdf) -- (chunk);
  \draw[arr] (chunk) -- (emb);
  \draw[arr] (emb) -- (llm);
  \draw[arr, dashed, red!70] (pdf) |- (fb);
  \draw[arr, dashed, red!70] (emb) -- (fb);
  \draw[arr, dashed, red!70] (fb) -| (llm);
\end{tikzpicture}}
\caption{RAG 笔记生成流程}
\end{figure}

\subsection{深度分析}
分析阶段生成三类结果：文献对比、研究脉络、研究空白。结果保存到会话的 \texttt{analysis} 目录中，核心文件为 \texttt{analysis\_results.json}，前端以三张 Markdown 卡片展示并允许编辑。

\subsection{Writer}
Writer 输入由三部分构成：聚合笔记、用户反馈、分析洞察。若分析文件不存在，则仅使用笔记，保证兼容。

%============================================================
\section{Skills 与可观测性}
%============================================================

\subsection{双通道注入}
\begin{figure}[H]
\centering
\resizetikz[0.88\textwidth]{\begin{tikzpicture}[
  node distance=1.1cm,
  box/.style={rectangle, draw=blue!65, fill=blue!6, rounded corners=2mm, text width=3.3cm, align=center, minimum height=0.8cm},
  decision/.style={diamond, draw=orange!80, fill=orange!8, aspect=2.3, text width=2.7cm, align=center},
  arr/.style={-{Stealth}, thick}
]
  \node[box] (phase) {阶段入口\\search / notes / write};
  \node[decision, right=of phase] (check) {Skill 是否有效};
  \node[box, above right=0.8cm and 1.2cm of check] (skill) {通道 A\\注入 Skill};
  \node[box, below right=0.8cm and 1.2cm of check] (default) {通道 B\\默认策略};
  \node[box, right=4.0cm of check] (trace) {写入\\SKILL\_STATUS};
  \draw[arr] (phase) -- (check);
  \draw[arr] (check) -- node[above]{是} (skill);
  \draw[arr] (check) -- node[below]{否} (default);
  \draw[arr] (skill) -- (trace);
  \draw[arr] (default) -- (trace);
\end{tikzpicture}}
\caption{Skill 双通道注入与可观测性}
\end{figure}

\subsection{Trace 字段}
\begin{lstlisting}[language=Python]
{
  "action": "SKILL_STATUS",
  "input": {
    "phase": "write",
    "skill_id": "skill_xxx",
    "skill_title": "custom_write_skill",
    "loaded": true,
    "fallback_default": false
  },
  "error_type": "skill_info"
}
\end{lstlisting}

%============================================================
\section{全局知识库与 Copilot}
%============================================================

\texttt{GlobalKnowledgeBase} 扫描所有真实 Session 的论文摘要、笔记和最新综述草稿，建立轻量检索索引。全局 Copilot 默认跨所有 Session 检索，也支持前端传入 \texttt{session\_ids} 限定范围。

\begin{figure}[H]
\centering
\resizetikz{\begin{tikzpicture}[
  node distance=1.3cm,
  store/.style={cylinder, shape border rotate=90, aspect=0.25, draw=green!55!black, fill=green!6, text width=2.4cm, align=center},
  box/.style={rectangle, draw=blue!70, fill=blue!6, rounded corners=2mm, text width=3.0cm, align=center},
  arr/.style={-{Stealth}, thick}
]
  \node[store] (s1) {Session A};
  \node[store, below=0.7cm of s1] (s2) {Session B};
  \node[store, below=0.7cm of s2] (s3) {Session C};
  \node[box, right=2.5cm of s2] (kb) {GlobalKnowledgeBase\\索引构建};
  \node[box, right=2.5cm of kb] (filter) {session\_ids\\范围过滤};
  \node[box, right=2.5cm of filter] (copilot) {全局 Copilot\\回答生成};
  \draw[arr] (s1) -- (kb);
  \draw[arr] (s2) -- (kb);
  \draw[arr] (s3) -- (kb);
  \draw[arr] (kb) -- (filter);
  \draw[arr] (filter) -- (copilot);
\end{tikzpicture}}
\caption{全局知识库检索范围控制}
\end{figure}

%============================================================
\section{测试与评估设计}
%============================================================

系统保留 \texttt{evaluation/} 目录，包含评估任务、数据库、脚本和说明。现有自动化测试通过 \texttt{pytest -q} 执行。后续评估应重点覆盖自动模式、分析注入、Skill trace 与 Copilot 范围选择。

\end{document}
